\documentclass[12pt]{article}
\usepackage{tgtermes}
\usepackage[T1]{fontenc}
\usepackage{lmodern}
\usepackage[margin=1in]{geometry}
\usepackage{setspace}
\usepackage{lineno}
\usepackage{booktabs}
\usepackage{amsmath}
\usepackage{longtable}
\usepackage{hyperref}

% -- breakable file paths: \path{} allows line breaks inside long
%    directory/file names, which \texttt{} cannot do --
\usepackage{url}
\urlstyle{tt}
\def\UrlBreaks{\do\/\do\_\do\-\do\.}

\doublespacing
\linenumbers

\title{Online Resource 1\\[0.5em]
\large Activation and Foreclosure: Separating Screening from Exit Costs\\in Codified Moral Communities}
\author{}
\date{}

\begin{document}
\maketitle

\noindent This supplement documents the parameterization of the simulation model,
the robustness of the outcome classification, the continuous surfaces that
underlie the discretized results, the privilege ablation, the endogenous
exit-capacity module, and the computational provenance of every reported
quantity. Section~S1 is taken from the frozen model sources. Sections~S2
through~S5 report quantities computed from the committed run outputs.
Section~S6 states where those outputs and the code that produced them reside.

\section{Parameter-to-observable mapping}
\label{sec:observable-mapping}

This section supplies the detailed measurement material moved from the main
text. The mapping proceeds from model parameters to institutional indicators.
Legibility is the extent to which compliance is classifiable through visible
acts, speech, dress, attendance, dietary practice, or other public markers.
Substitutability asks whether visible compliance is accepted as sufficient, or
near-sufficient, for moral standing. Enforcement affordance is the ease with
which deviations are detected, coded, and linked to sanctions. Centralization
is the concentration of interpretive and punitive authority. Exit cost is
approximated by legal penalties, kinship endogamy, economic dependence,
identity foreclosure, and the social cost of disaffiliation.

Enforcement reward is observed in promotion, prestige, institutional access,
or material compensation tied to disciplinary activity. Delegation is observed
in specialized offices, enforcement jurisdictions, and recognized roles for
monitoring and sanctioning. These indicators locate a bounded regime in a
common structural space; they do not produce a ranked taxonomy of traditions.

\begin{table}[htbp]
\centering
\caption{Code geometry dimensions and observable indicators.}
\begin{tabular}{lp{5.5cm}p{5.5cm}}
\toprule
Dimension & Model parameter & Observable indicators \\
\midrule
Legibility & $\sigma$ & Visible prescriptions, dress, attendance, dietary markers \\
Substitutability & Collapsed into $\sigma$ & Visible compliance accepted as sufficient for standing \\
Enforcement affordance & $\pi$, $\kappa$ & Violation categories, complaint channels, sanction pathways \\
Centralization & Cadre and monopoly & Adjudicating bodies, orthodoxy monopoly, restricted appeals \\
Exit capacity & base opportunity, $\delta$ & Legal, kinship, economic, identity, and social outside options \\
\bottomrule
\end{tabular}
\end{table}

\begin{table}[htbp]
\centering
\caption{Main-text model parameters and observable interpretation.}
\begin{tabular}{llll}
\toprule
Parameter & Symbol & Reported value & Interpretation \\
\midrule
Legibility weight & $\sigma$ & swept & Weight of external compliance \\
Enforcement reward & $\pi$ & swept & Return to successful punishment \\
Enforcement cost & $\kappa$ & 0.08 & Backlash risk per act \\
Enforcer quota & $q$ & 0.08 & Fraction designated as cadre \\
Alt. degradation & $\delta$ & swept & Loss of outside-option quality \\
Drift speed & $\eta$ & swept & Rate of endogenous degradation \\
\bottomrule
\end{tabular}
\end{table}

No single dimension is sufficient to determine regime type. High legibility
with open exit can produce active but leaky enforcement; high exit cost without
activation conditions can retain members without producing an apparatus.
The same tradition may occupy different locations across periods, states, or
rival institutions. The mapping compares configurations and does not measure
latent interior states.
\section{Model parameters}
\label{sec:params}

Two frozen model sources generate the reported results. The main experiments
(activation boundary, imposed-closure dose-response, shock factorial,
closure-decoupling, and privilege ablation) use
\path{src/religion_fundamentalism_abm_v2_7.py}. The endogenous exit-capacity
experiments of Section~S5 use \path{src/religion_fundamentalism_abm_v3_3.py},
which reproduces the former exactly when \texttt{exit\_capacity\_mode} is
\texttt{off}.

Tables~S1.1 through~S1.8 list every parameter of the main model with the
default value declared in the frozen source and the role that value plays.
The rightmost column records how the reported experiments set the parameter:
\emph{default} means the source default is used, \emph{swept} means the
experimental design varies it across the values given in the Methods, and
\emph{set} means the sweep launcher passes an explicit value that differs from
the source default. Table~S1.9 lists the additional parameters introduced by
the endogenous exit-capacity module. Table~S1.10 gives the calibration parameter
set that the sweep launcher supplies, which is provided in the repository at
\path{results/BEST_PARAMS_frozen.json}.

\begin{table}[htbp]
\centering
\footnotesize
\caption{Table S1.1: Population, network, and schedule.}
\begin{tabular}{p{4.9cm}p{1.3cm}p{4.5cm}p{2.4cm}}
\toprule
Parameter & Default & Role & In reported runs \\
\midrule
\texttt{n} & 300 & Number of agents at initialization & set to 350 \\
\texttt{graph} & scale\_free & Network generator; Barab\'asi--Albert with attachment parameter $m=\max\{2,\operatorname{round}[\log(N+1)]\}$ & default \\
\texttt{seed} & 1 & Seed for network generation, trait draws, and all stochastic updates & swept, 1 to 30 \\
\texttt{steps} & 350 & Number of simulated steps & set to 450 \\
\texttt{fixed\_y0} & False & When true, freezes the orthodox belief position at its initial value & default \\
\bottomrule
\end{tabular}
\end{table}

\begin{table}[htbp]
\centering
\footnotesize
\caption{Table S1.2: Code geometry and observation.}
\begin{tabular}{p{4.9cm}p{1.3cm}p{4.5cm}p{2.4cm}}
\toprule
Parameter & Default & Role & In reported runs \\
\midrule
\texttt{sigma} & 0.75 & Weight of externally legible compliance relative to intrinsic cultivation in the perceived-morality signal & swept \\
\texttt{v\_obs} & 0.90 & Probability that a compliance marker is observed & default \\
\texttt{a\_obs} & 0.05 & Probability that intrinsic cultivation is observed & default \\
\texttt{h\_obs} & 0.15 & Probability that an agent's belief position is observed for the heresy-distance term & default \\
\bottomrule
\end{tabular}
\end{table}

\begin{table}[htbp]
\centering
\footnotesize
\caption{Table S1.3: Agent behaviour, learning, and social dynamics.}
\begin{tabular}{p{4.9cm}p{1.3cm}p{4.5cm}p{2.4cm}}
\toprule
Parameter & Default & Role & In reported runs \\
\midrule
\texttt{c\_epc} & 0.12 & Utility cost of displaying the visible compliance marker & default \\
\texttt{eps\_cult} & 0.06 & Utility cost of investing in intrinsic cultivation & default \\
\texttt{eta\_cult} & 0.07 & Utility return to intrinsic cultivation & default \\
\texttt{beta\_epc} & 5.0 & Logistic slope converting net status gain into the probability of displaying compliance & default \\
\texttt{beta\_pun} & 6.0 & Logistic slope converting the perceived compliance deficit into punishment probability & default \\
\texttt{k\_h} & 10.0 & Logistic slope of the heresy-distance term & default \\
\texttt{rho\_r} & 0.04 & Step size for updating enforcement propensity after a punishment outcome & default \\
\texttt{rho\_b} & 0.03 & Step size for updating rigidity after a punishment outcome & default \\
\texttt{w\_m} & 1.0 & Weight on the perceived-morality shortfall in the punishment deficit & default \\
\texttt{w\_heresy} & 0.6 & Weight on the heresy-distance term in the punishment deficit & default \\
\texttt{d0} & 0.18 & Baseline heresy-distance threshold above which belief counts as deviant & default \\
\texttt{norm\_weight} & 0.55 & Weight on the local compliance norm in the status gain from displaying compliance & default \\
\texttt{shame\_weight} & 0.70 & Weight on the local enforcement norm in that status gain & default \\
\texttt{noise\_x} & 0.02 & Standard deviation of per-step noise on intrinsic cultivation & default \\
\texttt{cult\_step} & 0.04 & Maximum per-step change in intrinsic cultivation & default \\
\texttt{F\_star} & 0.45 & Threshold on the fundamentalism index used to report prevalence & default \\
\texttt{L\_alpha} & 1.5 & First shape parameter of the Beta distribution of literalism & default \\
\texttt{L\_beta} & 4.5 & Second shape parameter of that distribution & default \\
\texttt{theta\_L\_gain} & 0.12 & Sensitivity of an agent's effective strictness threshold to its literalism & default \\
\texttt{d0\_L\_gain} & 0.08 & Sensitivity of an agent's effective heresy threshold to its literalism & default \\
\bottomrule
\end{tabular}
\end{table}

\begin{table}[htbp]
\centering
\footnotesize
\caption{Table S1.4: Policing economy.}
\begin{tabular}{p{4.9cm}p{1.3cm}p{4.5cm}p{2.4cm}}
\toprule
Parameter & Default & Role & In reported runs \\
\midrule
\texttt{pi\_reward} & 0.22 & Coefficient on the legibility index in the institutional return to a successful punishment & swept \\
\texttt{kappa\_cost} & 0.08 & Base utility cost of issuing a punishment & default \\
\texttt{lam\_punish} & 0.25 & Utility penalty imposed on the punished agent & default \\
\texttt{backlash\_base} & 0.25 & Baseline probability that a punishment provokes backlash against the punisher & default \\
\texttt{backlash\_cost} & 0.30 & Utility cost borne by the punisher when backlash occurs & default \\
\texttt{backlash\_sensitivity} & 1.0 & Sensitivity of backlash probability to the illegibility of the code & default \\
\texttt{punish\_floor} & 0.08 & Minimum step-level fraction of active agents punished that gates any update of outside-option closure & set to 0.08 \\
\bottomrule
\end{tabular}
\end{table}

\begin{table}[htbp]
\centering
\footnotesize
\caption{Table S1.5: Exit desire, opportunity, and cost.}
\begin{tabular}{p{4.9cm}p{1.3cm}p{4.5cm}p{2.4cm}}
\toprule
Parameter & Default & Role & In reported runs \\
\midrule
\texttt{enable\_exit} & True & Master switch for the exit process & default \\
\texttt{exit\_threshold} & $-1.0$ & Utility level below which exit desire begins to accumulate & set to $-1.0$ \\
\texttt{exit\_prob\_slope} & 3.0 & Logistic slope converting exit drive into exit desire & default \\
\texttt{exit\_cost} & 0.40 & Base cost of leaving, before degree and lock-in adjustments & set to 1.00 \\
\texttt{exit\_degree\_coeff} & 0.03 & Rate at which network degree raises the effective exit cost & default \\
\texttt{exit\_epc\_lockin\_coeff} & 0.25 & Rate at which an agent's accumulated visible compliance raises its exit cost & default \\
\texttt{exit\_min\_cost} & 0.0 & Floor on the effective exit cost & default \\
\texttt{exit\_opportunity\_base} & 0.6 & Baseline logit intercept of the exit-opportunity draw & set to 0.30 \\
\texttt{exit\_opportunity\_deg\_coeff} & 0.03 & Rate at which degree lowers exit opportunity & default \\
\texttt{exit\_opportunity\_threat\_coeff} & 1.5 & Rate at which external threat lowers exit opportunity & set to 2.9081 \\
\texttt{exit\_block\_exponent} & 2.5 & Exponent applied to the raw opportunity, concentrating opportunity at high values & set to 5.9806 \\
\texttt{exit\_block\_floor} & 0.02 & Additive floor on exit opportunity after the exponent is applied & default \\
\texttt{exit\_commit\_steps} & 8 & Consecutive steps for which desire must persist before an opportunity draw is taken & set to 12 \\
\texttt{exit\_cooldown} & 0 & Steps an agent must wait after a failed exit attempt & default \\
\texttt{exit\_rewire\_fraction} & 0.90 & Fraction of an exiting agent's incident edges removed on departure & default \\
\texttt{alpha\_punish\_revalue} & 0.0 & Weight by which received punishment revalues staying inside the group & set to 0.00 \\
\texttt{mu\_membership\_reward} & 0.0 & Weight on the membership reward that offsets exit desire & set to 0.00 \\
\texttt{membership\_benefit} & 0.03 & Base per-step benefit of remaining a member & set to 0.0634 \\
\texttt{membership\_benefit\_sigma} & 0.08 & Sensitivity of that benefit to code legibility & default \\
\texttt{membership\_benefit\_threat} & 0.10 & Sensitivity of that benefit to external threat & set to 0.1055 \\
\bottomrule
\end{tabular}
\end{table}

\begin{table}[htbp]
\centering
\footnotesize
\caption{Table S1.6: Outside-option closure.}
\begin{tabular}{p{4.9cm}p{1.3cm}p{4.5cm}p{2.4cm}}
\toprule
Parameter & Default & Role & In reported runs \\
\midrule
\texttt{delta\_outside\_degrade} & 0.0 & Imposed baseline degradation of the outside option, denoted $\delta_0$ & swept \\
\texttt{eta\_delta\_drift} & 0.0 & Rate at which closure drifts toward its target, denoted $\eta$ & swept \\
\texttt{delta\_mode} & legacy & Selects the closure specification: \texttt{legacy} makes the target a function of enforcer share, \texttt{decoupled} makes it a function of active-agent punishment incidence, and \texttt{exogenous} holds closure fixed & swept \\
\texttt{delta\_cap} & 0.85 & Upper bound on the closure target in decoupled mode, denoted $\delta_{\max}$ & default \\
\texttt{delta\_kappa} & 3.0 & Coupling strength from punishment incidence to closure in decoupled mode & swept over 1.5, 3.0, 6.0 \\
\bottomrule
\end{tabular}
\end{table}

\begin{table}[htbp]
\centering
\footnotesize
\caption{Table S1.7: Institutional privilege bundle.}
\begin{tabular}{p{4.9cm}p{1.3cm}p{4.5cm}p{2.4cm}}
\toprule
Parameter & Default & Role & In reported runs \\
\midrule
\texttt{enforcer\_quota\_frac} & 0.08 & Target cadre size as a fraction of active agents & default \\
\texttt{cap\_to\_enforcer} & 0.25 & Institutional capital required for first-round cadre eligibility & default \\
\texttt{A\_enforcer\_monopoly} & 0.35 & Controller authority above which the punishment monopoly excludes non-enforcers & default \\
\texttt{enforcer\_punish\_mult} & 1.5 & Multiplier on an enforcer's punishment probability & default \\
\texttt{non\_enforcer\_punish\_eps} & 0.02 & Operative multiplier applied to non-enforcers when the monopoly is off & default \\
\texttt{non\_enforcer\_punish\_mult} & 0.25 & Declared non-enforcer multiplier; passed by the sweeps but not used in the operative branch & default \\
\texttt{kappa\_cap\_discount} & 0.20 & Fraction by which accumulated capital reduces the base punishment cost & default \\
\texttt{enforcer\_kappa\_mult} & 0.30 & Further multiplier on an enforcer's punishment cost, a 70 percent role discount & default \\
\texttt{enforcer\_backlash\_mult} & 0.25 & Multiplier on an enforcer's backlash probability, 75 percent protection & default \\
\texttt{cap\_gain\_per\_punish} & 0.15 & Institutional capital added per successful punishment & default \\
\texttt{cap\_max} & 2.0 & Ceiling on accumulated institutional capital & default \\
\texttt{cap\_decay} & 0.005 & Per-step decay rate of institutional capital & default \\
\texttt{service\_decay} & 0.02 & Per-step decay rate of accumulated service & default \\
\texttt{budget\_base} & 0.15 & Constant term of the controller patronage budget & default \\
\texttt{budget\_threat\_gain} & 0.60 & Coefficient on threat in the patronage budget & default \\
\bottomrule
\end{tabular}
\end{table}

\begin{table}[htbp]
\centering
\footnotesize
\caption{Table S1.8: Institutional controller and threat schedule.}
\begin{tabular}{p{4.9cm}p{1.3cm}p{4.5cm}p{2.4cm}}
\toprule
Parameter & Default & Role & In reported runs \\
\midrule
\texttt{shock\_schedule} & 100, 220, 320 & Steps at which external threat is raised & swept: empty, or 100, 220, 320 \\
\texttt{shock\_strength} & 0.25 & Increment added to threat at each scheduled shock & set to 0.10 \\
\texttt{relax\_rate} & 0.03 & Per-step proportional decay of threat between shocks & default \\
\texttt{A\_gain\_threat} & 2.0 & Rate at which threat raises controller authority through $A_t=1-\exp(-A_{\mathrm{gain}}T_t)$ & default \\
\texttt{tighten\_gain\_pi} & 0.20 & Increase in enforcement reward per unit threat & default \\
\texttt{tighten\_gain\_lam} & 0.20 & Increase in punishment severity per unit threat & default \\
\texttt{tighten\_gain\_d0} & 0.08 & Reduction in the heresy threshold per unit threat & default \\
\texttt{baseline\_pi} & $-1.0$ & Sentinel; a negative value makes the run adopt \texttt{pi\_reward} as the baseline & default \\
\texttt{baseline\_lam} & $-1.0$ & Sentinel for the punishment-severity baseline & default \\
\texttt{baseline\_d0} & $-1.0$ & Sentinel for the heresy-threshold baseline & default \\
\bottomrule
\end{tabular}
\end{table}

\paragraph{Note on launcher-supplied values.}
The sweep launcher \path{scripts/run_v2_7_endogenous_delta_sweep.py} reads seven
parameters from a calibration file supplied through its \texttt{-{}-best}
argument: \path{exit_cost}, \path{exit_block_exponent},
\path{exit_commit_steps}, \path{exit_opportunity_threat_coeff},
\path{membership_benefit}, \path{membership_benefit_threat}, and
\path{shock_strength}. Its \texttt{load\_best\_params} function reads these seven
keys and its \texttt{run\_one} function passes each of them to the model as a
command-line flag. For all seven, the value that ran differs from the frozen
source default, so each is a \emph{set} row above rather than a \emph{default}
row.

The calibration file is provided in the repository at
\path{results/BEST_PARAMS_frozen.json}. Table~S1.10 gives its contents.

\begin{table}[htbp]
\centering
\footnotesize
\caption{Table S1.10: The calibration parameter set, in
\texttt{results/BEST\_PARAMS\_frozen.json}.}
\begin{tabular}{p{5.4cm}p{2.2cm}p{2.2cm}p{3.4cm}}
\toprule
Parameter & Source default & Value that ran & How it was confirmed \\
\midrule
\texttt{exit\_cost} & 0.40 & 1.0 & algebraic inversion; replay \\
\texttt{exit\_block\_exponent} & 2.5 & 5.980606901200484 & replay \\
\texttt{exit\_commit\_steps} & 8 & 12 & disengagement threshold; replay \\
\texttt{exit\_opportunity\_threat\_coeff} & 1.5 & 2.908096432532684 & replay, shock-on only \\
\texttt{membership\_benefit} & 0.03 & 0.06336634914903713 & \texttt{stay\_value} column; replay \\
\texttt{membership\_benefit\_threat} & 0.10 & 0.10553999169992101 & \texttt{stay\_value} regression, shock-on only \\
\texttt{shock\_strength} & 0.25 & 0.1 & \texttt{threat} column, shock-on only \\
\bottomrule
\end{tabular}
\end{table}

An eighth key, \texttt{exit\_opportunity\_base}, is present in the calibration
file but is not read by \texttt{load\_best\_params}. The launcher passes its own
\texttt{-{}-base-opp}, which is 0.30 in every reported sweep and is recorded as
\texttt{base\_opp=0.3} in each \path{sweep_report.txt}.

\paragraph{How the seven values were confirmed.}
The committed run outputs do not record a command line, so each value was
recovered from what the runs wrote and then checked by replay. Three methods
agree.

First, algebraic inversion. \path{agent_summary.csv} records
\texttt{exit\_cost\_eff\_last}, which the model computes as
$\texttt{exit\_cost}\,(1+0.03d_i)+0.25s_i$ with $s_i\in\{0,1\}$. Inverting this
over every agent with a nonzero entry returns exactly 1.0 and never 0.40.

Second, direct reading of echoed columns. In \path{recon/factorial_shock_on}, the
\texttt{threat} column steps from 0 to exactly 0.1 across the first scheduled
shock and peaks at 0.1050330, the value implied by an increment of 0.1 under the
0.97 per-step relaxation. Regressing the \texttt{stay\_value} column on
\texttt{threat} in the same run returns an intercept of 0.0843663 and a slope of
0.1055400, which give \texttt{membership\_benefit} $=0.0633663$ and
\texttt{membership\_benefit\_threat} $=0.1055400$ once the legibility term is
subtracted. The agent-level \texttt{disengaged} flag is set at
$\lfloor\texttt{exit\_commit\_steps}/2\rfloor$; its observed threshold of 6
implies 12 and excludes 8.

Third, replay, which is decisive. Re-running
\path{src/religion_fundamentalism_abm_v2_7.py} at the recorded cell and seed
reproduces the committed \path{metrics.csv} and \path{agent_summary.csv} exactly
under the parameter set in Table~S1.10, for one run drawn from each of
\path{recon/exogenous_delta_fixed}, \path{recon/boundary_sealed}, and
\path{recon/factorial_shock_on}. Substituting the frozen-source defaults breaks
the reproduction in all three: terminal exit moves from 0.046 to 0.871 in the
first run and from 0.003 to 0.354 in the second.

\paragraph{Identifiability caveat.}
The three parameters \path{shock_strength},
\path{exit_opportunity_threat_coeff}, and
\path{membership_benefit_threat} all enter the model multiplied by threat.
Every sweep except \path{recon/factorial_shock_on} runs with an empty shock
schedule, so threat is identically zero and those three parameters are inert and
formally unidentifiable there. This was confirmed directly: replaying the
\path{recon/exogenous_delta_fixed} run with those three set to the source
defaults still reproduces the committed output bit-for-bit. They are identified
only in \path{recon/factorial_shock_on}, where setting
\texttt{exit\_opportunity\_threat\_coeff} to 1.5 breaks bit-for-bit identity and
2.908096432532684 restores it.

\begin{table}[htbp]
\centering
\footnotesize
\caption{Table S1.9: Additional parameters of the endogenous exit-capacity module.}
\begin{tabular}{p{4.9cm}p{1.3cm}p{4.5cm}p{2.4cm}}
\toprule
Parameter & Default & Role & In reported runs \\
\midrule
\texttt{exit\_capacity\_mode} & off & Selects the capacity specification; \texttt{off} reproduces the main model exactly, \texttt{endogenous} replaces global closure with a per-agent capacity & set to endogenous \\
\texttt{ec\_init\_random} & True & Draws per-agent outside ties and economic independence from Beta(2,2) at initialization & default \\
\texttt{ec\_outside\_ties\_init} & 1.0 & Initial outside-tie level when homogeneous initialization is selected & default \\
\texttt{ec\_tie\_decay} & 0.01 & Per-step decay rate of outside social ties & default \\
\texttt{ec\_tie\_renewal\_base} & 0.006 & Constant term of outside-tie renewal; the tie channel closes when crowd-out exceeds this value & default \\
\texttt{ec\_tie\_renewal} & 0.004 & Renewal of outside ties proportional to current ties & default \\
\texttt{ec\_tenure\_crowdout} & 0.002 & Rate at which tenure inside the group crowds out outside ties & swept over 0.0005, 0.002, 0.008 \\
\texttt{ec\_econ\_recovery\_base} & 0.004 & Constant term of economic-independence recovery & default \\
\texttt{ec\_econ\_recovery} & 0.002 & Recovery of economic independence proportional to its current level & default \\
\texttt{ec\_dependence\_rate} & 0.01 & Rate at which membership reward erodes economic independence & default \\
\texttt{ec\_hetero\_sd} & 0.25 & Standard deviation of per-agent heterogeneity in capacity & default \\
\texttt{lambda\_exit\_opportunity} & 1.0 & Weight of capacity on the exit-opportunity channel & default \\
\texttt{lambda\_exit\_willingness} & 1.0 & Weight of capacity on the exit-willingness channel & default \\
\texttt{turnover\_mode} & off & Enables cohort turnover through death and replacement & set to off \\
\texttt{death\_hazard} & 0.002 & Per-step death hazard when turnover is on & default \\
\texttt{born\_inside\_frac} & 0.8 & Fraction of the population initialized as born inside the group & swept over 0.0 and 0.8 \\
\texttt{emit\_panel} & False & Emits the per-step per-agent panel used for the anti-circularity test & set to on for the panel runs \\
\bottomrule
\end{tabular}
\end{table}

\clearpage

\section{Classification robustness}
\label{sec:robustness}

The canonical classifier defines capture as retention together with activation:
a run is retained when \texttt{final\_exit\_rate} $\leq 0.20$ and active when
\texttt{max\_punish} $\geq 0.10$. This section reports the sensitivity of that
definition to the activation threshold, the vacuity of the concentration
criterion that the definition omits, and the relationship between COLLAPSE and
activation.

\subsection{Activation-threshold sensitivity}

Table~S2.1 varies the activation threshold and records, for each value, the
capture rate in the imposed-closure experiment, the capture rate at the sealed
boundary, whether the imposed-closure dose-response remains monotone, and
whether the low-observability, low-reward cell at $\sigma=0.25$, $\pi=0.05$
remains inactive.

\begin{table}[htbp]
\centering
\footnotesize
\caption{Table S2.1: Capture rate and qualitative behaviour against the activation threshold.}
\begin{tabular}{p{2.0cm}p{2.5cm}p{2.5cm}p{2.2cm}p{3.6cm}}
\toprule
Threshold & Imposed-closure capture & Sealed-boundary capture & Dose-response & $\sigma=0.25$, $\pi=0.05$ cell \\
\midrule
0.05 & 0.253 & 0.816 & monotone & breaks, 90 of 180 capture \\
0.075 & 0.198 & 0.694 & monotone & holds, 0 of 180 \\
0.10 (canonical) & 0.198 & 0.661 & monotone & holds, 0 of 180 \\
0.15 & 0.186 & 0.473 & monotone & holds, 0 of 180 \\
0.20 & 0.006 & 0.070 & degenerates & holds, 0 of 180 \\
\bottomrule
\end{tabular}
\end{table}

Both the dose-response and the necessity result survive across the band
$[0.075, 0.15]$, and the canonical value 0.10 lies in the interior of that
band. The two edges fail in opposite directions. At the loose edge, 0.05, the
$\sigma=0.25$, $\pi=0.05$ cell ceases to be inactive and returns 90 of 180 runs
as capture, so the necessity result no longer holds. At the strict edge, 0.20,
capture falls to approximately zero everywhere, with an imposed-closure capture
rate of 0.006 and a sealed-boundary rate of 0.070, so the dose-response
degenerates and carries no information. The reported band is therefore joint:
it is the set of thresholds at which both results hold simultaneously. Sources:
\path{recon/exogenous_delta_fixed} and \path{recon/boundary_sealed}.

\subsection{Vacuity of the concentration criterion}

An earlier form of the classifier gated capture on an enforcer punishment share
of at least 0.70. Reclassifying every committed run with and without that gate
changes 0 of 9,570 classifications. The 9,570 runs comprise 4,050 from the
shock-off factorial, 1,620 from the imposed-closure experiment, 1,890 from the
sealed boundary, 1,890 from the open boundary, and 120 from the ablation
ceiling.

The reason the gate is inert is distributional. Enforcer share has no mass near
zero in any of the five datasets: the fraction of runs at or below 0.05 is 0.000
throughout, and the median lies between 0.97 and 0.98. The statistic saturates
high wherever punishment occurs at all, which makes it redundant with activation
rather than a graded measure of concentration. It is therefore reported
descriptively and excluded from the capture definition.

\subsection{COLLAPSE and activation}

Every COLLAPSE run in the committed data falls below the canonical activation
threshold. Table~S2.2 gives the count and the observed range of maximum
punishment incidence by dataset.

\begin{table}[htbp]
\centering
\footnotesize
\caption{Table S2.2: COLLAPSE runs and their maximum punishment incidence.}
\begin{tabular}{p{6.0cm}p{1.9cm}p{2.9cm}p{2.6cm}}
\toprule
Dataset & COLLAPSE runs & Of which \texttt{max\_punish} $<0.10$ & min / median / max \\
\midrule
\path{results/v2.5_methodology_paper_canonical} & 45 & 45 (100 percent) & 0.037 / 0.063 / 0.094 \\
\path{results/v2.5_corrected_three_regime_confirm} & 45 & 45 (100 percent) & 0.037 / 0.063 / 0.094 \\
\path{results/v2.5_corrected_sweep_regime_search_fast} & 18 & 18 (100 percent) & 0.037 / 0.059 / 0.094 \\
All \path{recon/} sweeps & 0 & not applicable & not applicable \\
\bottomrule
\end{tabular}
\end{table}

All 108 COLLAPSE runs are inactive, and the largest maximum punishment incidence
observed among them is 0.094, below the 0.10 activation threshold. Depopulation
in these runs is therefore driven by open exit rather than by enforcement, and
COLLAPSE is properly described as depopulation without activation.

This has a direct consequence for the structure of the classifier. Reclassifying
all committed data as activation crossed with retention, with no separate
COLLAPSE rule, changes 108 of 28,581 runs, or 0.38 percent, and every change is
COLLAPSE to QUIET. The inactive cell with exit strictly between 0.20 and 0.90
contains 560 runs, or 1.96 percent, all already labelled QUIET. COLLAPSE is
retained only as a descriptive flag for the exit $\geq 0.90$ tail within the
inactive region, not as a separate dynamical regime.

\clearpage

\section{Continuous surfaces}
\label{sec:continuous}

The classification in Section~S2 discretizes two continuous quantities. This
section reports those quantities directly, so that the reported results can be
read without reference to any threshold. Each entry is a mean over 30 seeds for
the corresponding cell.

\subsection{Open exit}

The open-exit condition fixes $\delta_0=0$ and comprises 1,890 runs on a grid of
9 values of $\sigma$ by 7 values of $\pi$ by 30 seeds. Source:
\path{recon/boundary_open/sweep_seed_results.csv}.

\begin{table}[htbp]
\centering
\small
\caption{Table S3.1: Mean terminal exit rate, open exit.}
\begin{tabular}{lccccccc}
\toprule
$\sigma \backslash \pi$ & 0.01 & 0.03 & 0.05 & 0.10 & 0.15 & 0.25 & 0.50 \\
\midrule
0.05 & 0.176 & 0.173 & 0.176 & 0.174 & 0.172 & 0.542 & 0.500 \\
0.15 & 0.096 & 0.097 & 0.094 & 0.374 & 0.490 & 0.464 & 0.440 \\
0.25 & 0.046 & 0.044 & 0.044 & 0.446 & 0.433 & 0.405 & 0.426 \\
0.35 & 0.016 & 0.016 & 0.319 & 0.409 & 0.383 & 0.403 & 0.446 \\
0.45 & 0.005 & 0.005 & 0.400 & 0.385 & 0.382 & 0.400 & 0.448 \\
0.55 & 0.002 & 0.219 & 0.380 & 0.384 & 0.393 & 0.417 & 0.440 \\
0.65 & 0.000 & 0.306 & 0.368 & 0.378 & 0.384 & 0.424 & 0.419 \\
0.75 & 0.000 & 0.343 & 0.366 & 0.374 & 0.389 & 0.400 & 0.399 \\
0.95 & 0.000 & 0.319 & 0.318 & 0.324 & 0.322 & 0.317 & 0.310 \\
\bottomrule
\end{tabular}
\end{table}

\begin{table}[htbp]
\centering
\small
\caption{Table S3.2: Mean maximum punishment incidence, open exit.}
\begin{tabular}{lccccccc}
\toprule
$\sigma \backslash \pi$ & 0.01 & 0.03 & 0.05 & 0.10 & 0.15 & 0.25 & 0.50 \\
\midrule
0.05 & 0.050 & 0.050 & 0.050 & 0.050 & 0.050 & 0.091 & 0.110 \\
0.15 & 0.050 & 0.050 & 0.050 & 0.064 & 0.098 & 0.109 & 0.148 \\
0.25 & 0.051 & 0.051 & 0.051 & 0.102 & 0.110 & 0.132 & 0.172 \\
0.35 & 0.046 & 0.047 & 0.074 & 0.117 & 0.130 & 0.156 & 0.181 \\
0.45 & 0.046 & 0.046 & 0.105 & 0.133 & 0.152 & 0.173 & 0.181 \\
0.55 & 0.044 & 0.066 & 0.126 & 0.149 & 0.164 & 0.173 & 0.180 \\
0.65 & 0.042 & 0.093 & 0.136 & 0.156 & 0.163 & 0.171 & 0.174 \\
0.75 & 0.039 & 0.118 & 0.143 & 0.154 & 0.162 & 0.162 & 0.163 \\
0.95 & 0.033 & 0.128 & 0.130 & 0.136 & 0.138 & 0.136 & 0.135 \\
\bottomrule
\end{tabular}
\end{table}

\subsection{Sealed exit}

The sealed-exit condition fixes $\delta_0=0.95$ on the same grid, again 1,890
runs. Source: \path{recon/boundary_sealed/sweep_seed_results.csv}.

\begin{table}[htbp]
\centering
\small
\caption{Table S3.3: Mean terminal exit rate, sealed exit.}
\begin{tabular}{lccccccc}
\toprule
$\sigma \backslash \pi$ & 0.01 & 0.03 & 0.05 & 0.10 & 0.15 & 0.25 & 0.50 \\
\midrule
0.05 & 0.008 & 0.008 & 0.009 & 0.007 & 0.008 & 0.084 & 0.107 \\
0.15 & 0.004 & 0.004 & 0.004 & 0.047 & 0.080 & 0.099 & 0.138 \\
0.25 & 0.002 & 0.002 & 0.002 & 0.077 & 0.081 & 0.105 & 0.141 \\
0.35 & 0.001 & 0.001 & 0.042 & 0.076 & 0.096 & 0.128 & 0.128 \\
0.45 & 0.000 & 0.000 & 0.068 & 0.089 & 0.110 & 0.117 & 0.118 \\
0.55 & 0.000 & 0.023 & 0.068 & 0.100 & 0.108 & 0.113 & 0.110 \\
0.65 & 0.000 & 0.041 & 0.073 & 0.092 & 0.100 & 0.097 & 0.097 \\
0.75 & 0.000 & 0.058 & 0.076 & 0.081 & 0.083 & 0.082 & 0.083 \\
0.95 & 0.000 & 0.056 & 0.053 & 0.055 & 0.055 & 0.054 & 0.054 \\
\bottomrule
\end{tabular}
\end{table}

\begin{table}[htbp]
\centering
\small
\caption{Table S3.4: Mean maximum punishment incidence, sealed exit.}
\begin{tabular}{lccccccc}
\toprule
$\sigma \backslash \pi$ & 0.01 & 0.03 & 0.05 & 0.10 & 0.15 & 0.25 & 0.50 \\
\midrule
0.05 & 0.050 & 0.050 & 0.050 & 0.050 & 0.050 & 0.107 & 0.126 \\
0.15 & 0.050 & 0.050 & 0.050 & 0.087 & 0.113 & 0.131 & 0.186 \\
0.25 & 0.051 & 0.051 & 0.051 & 0.121 & 0.129 & 0.164 & 0.196 \\
0.35 & 0.046 & 0.047 & 0.099 & 0.136 & 0.159 & 0.192 & 0.202 \\
0.45 & 0.046 & 0.046 & 0.130 & 0.163 & 0.183 & 0.197 & 0.204 \\
0.55 & 0.044 & 0.086 & 0.147 & 0.182 & 0.192 & 0.200 & 0.202 \\
0.65 & 0.042 & 0.125 & 0.161 & 0.179 & 0.189 & 0.197 & 0.195 \\
0.75 & 0.039 & 0.151 & 0.173 & 0.179 & 0.189 & 0.189 & 0.189 \\
0.95 & 0.033 & 0.154 & 0.158 & 0.164 & 0.166 & 0.165 & 0.165 \\
\bottomrule
\end{tabular}
\end{table}

Comparing Tables~S3.1 and~S3.3 shows that sealing the exit lowers terminal exit
throughout the grid. Comparing Tables~S3.2 and~S3.4 shows that it also raises
enforcement intensity across most of the grid, which is the continuous form of
the coupling between the two axes at the frontier margin.

\subsection{Imposed closure}

The imposed-closure experiment holds closure fixed at a series of levels and
comprises 1,620 runs over 3 values of $\sigma$ by 3 values of $\pi$ by 6 values
of $\delta_0$ by 30 seeds. Source:
\path{recon/exogenous_delta_fixed/sweep_seed_results.csv}.

\begin{table}[htbp]
\centering
\small
\caption{Table S3.5: Mean maximum punishment and mean terminal exit by imposed closure, pooled over $\sigma$ and $\pi$.}
\begin{tabular}{lccc}
\toprule
$\delta_0$ & Mean maximum punishment & Mean terminal exit & $n$ \\
\midrule
0.00 & 0.136 & 0.332 & 270 \\
0.20 & 0.139 & 0.324 & 270 \\
0.40 & 0.143 & 0.313 & 270 \\
0.60 & 0.148 & 0.288 & 270 \\
0.80 & 0.155 & 0.216 & 270 \\
0.95 & 0.161 & 0.072 & 270 \\
\bottomrule
\end{tabular}
\end{table}

\begin{table}[htbp]
\centering
\small
\caption{Table S3.6: Mean maximum punishment over $\sigma$ and $\pi$ at $\delta_0=0.95$.}
\begin{tabular}{lccc}
\toprule
$\sigma \backslash \pi$ & 0.05 & 0.25 & 0.50 \\
\midrule
0.25 & 0.051 & 0.164 & 0.196 \\
0.75 & 0.173 & 0.189 & 0.189 \\
0.95 & 0.158 & 0.165 & 0.165 \\
\bottomrule
\end{tabular}
\end{table}

Raising imposed closure from 0.00 to 0.95 lowers mean terminal exit from 0.332
to 0.072 while raising mean maximum punishment only from 0.136 to 0.161. Closure
acts principally on retention rather than on activation.

\subsection{The necessity result without a threshold}

The necessity claim has a form that requires no classification. At the
low-observability, low-reward corner $\sigma=0.25$, $\pi=0.05$, mean maximum
punishment over 30 seeds is 0.051 at $\delta_0=0.0$ and 0.051 at
$\delta_0=0.95$, a change of $+0.000$, while mean terminal exit falls from 0.044
to 0.002. Sealing the exit in this corner removes the departures without
producing any enforcement. Both values lie far below the 0.10 activation band,
so the statement is a claim about the continuous enforcement surface and does
not depend on where the activation threshold is placed.

\clearpage

\section{Privilege ablation}
\label{sec:ablation}

The ablation isolates the contribution of the institutional privilege bundle to
the concentration of punishment. It compares a floor arm, in which every
privilege is removed, against a ceiling arm carrying the complete bundle, and
then adds each privilege back individually. All metrics are computed over agents
active at the end of the run, and each arm contains 120 runs. Source:
\path{recon/privilege_ablation/ablation_seed_results.csv}.

\subsection{Floor, ceiling, and the random-allocation null}

The null holds the observed active-agent punishment volume fixed and reallocates
those events uniformly at random across active agents, which gives the
concentration that arises from counting alone.

\begin{table}[htbp]
\centering
\small
\caption{Table S4.1: Floor and ceiling concentration against the random-allocation null.}
\begin{tabular}{lcccc}
\toprule
Arm & Top-5 percent share & Gini & Gini, random null & Terminal exit rate \\
\midrule
Floor & 0.091 & 0.240 & 0.028 & 0.585 \\
Ceiling & 0.580 & 0.901 & 0.062 & 0.331 \\
\bottomrule
\end{tabular}
\end{table}

The privilege-manufactured share of top-5 percent concentration is
$(0.580 - 0.091)/0.580 = 0.843$, that is 84 percent. The floor Gini of 0.240 sits
only modestly above its random null of 0.028, so the residual concentration in a
privilege-free model is small. Concentration is a product of the privilege
architecture rather than an emergent property of the underlying interaction.

\subsection{Single-privilege add-back}

Each privilege is added individually to the floor, holding all others removed,
with 120 seeds per arm. Recovery is expressed as a fraction of the floor to
ceiling span in top-5 percent share, which runs from 0.0907 to 0.5799.

\begin{table}[htbp]
\centering
\small
\caption{Table S4.2: Concentration after adding one privilege to the floor.}
\begin{tabular}{lccc}
\toprule
Arm & Top-5 percent share & Gini & Share of span recovered \\
\midrule
Floor & 0.0907 & 0.2402 & not applicable \\
Quota & 0.0907 & 0.2402 & 0.0 percent \\
Monopoly & 0.0907 & 0.2402 & 0.0 percent \\
Punishment multiplier & 0.2505 & 0.6317 & 32.7 percent \\
Backlash protection & 0.0907 & 0.2402 & 0.0 percent \\
Capital gain & 0.0907 & 0.2402 & 0.0 percent \\
Budget patronage & 0.0907 & 0.2402 & 0.0 percent \\
Cost discount & 0.0907 & 0.2402 & 0.0 percent \\
Ceiling & 0.5799 & 0.9007 & 100 percent \\
\bottomrule
\end{tabular}
\end{table}

Seven of the eight single privileges are exactly inert, reproducing the floor
values to four decimal places. The exception is the punishment multiplier, which
on its own recovers 32.7 percent of the span. The accurate statement is
therefore conjunctivity with one exception rather than conjunctivity outright:
concentration requires the privileges jointly, except for the coercive-power
multiplier, which is the single load-bearing privilege. The seven inert
privileges are inert for a common reason. Each acts only on a cadre, and the
quota-less floor never forms one, so there is no differentiated group for them
to act upon.

The multiplier exception carries a volume caveat that the concentration
statistics alone do not show. Table~S4.2 reports shares and Gini coefficients,
which are scale-free and can be computed over any nonzero punishment count. The
multiplier arm issues a mean of 393.2 punishments per run, against 75{,}611.8 at
the floor and 19{,}100.6 at the ceiling, a volume roughly two orders of
magnitude below either reference arm. It also retains almost the entire
population, with a mean terminal exit rate of 0.010 against 0.585 at the floor
and 0.331 at the ceiling, so the 32.7 percent recovery is a concentration
measured over about 1.1 punishment events per surviving active agent across 450
steps. The exception is real but thin: it establishes that the coercive-power
multiplier is the one privilege that can concentrate sanctioning without a
cadre, not that it reproduces the enforcement volume of the full bundle.

\subsection{Concentration does not discriminate capture}

Concentration is high wherever punishment occurs and therefore does not separate
the outcome classes. Table~S4.3 reports active-agent concentration by recomputed
regime in the imposed-closure experiment.

\begin{table}[htbp]
\centering
\small
\caption{Table S4.3: Concentration and exit by regime, imposed-closure experiment.}
\begin{tabular}{lcccc}
\toprule
Regime & $n$ & Top-5 percent share & Gini & Terminal exit rate \\
\midrule
CAPTURE & 320 & 0.678 & 0.918 & 0.105 \\
MIXED & 1120 & 0.625 & 0.903 & 0.339 \\
QUIET & 180 & 0.425 & 0.713 & 0.021 \\
\bottomrule
\end{tabular}
\end{table}

Among active systems, that is CAPTURE against MIXED, the two concentration
measures are nearly identical, 0.678 against 0.625 for top-5 percent share and
0.918 against 0.903 for the Gini, while terminal exit differs by roughly a factor
of three, 0.105 against 0.339. Retention rather than concentration separates the
regimes. QUIET sits lower on both concentration measures only because it is
barely active, which is an activation effect rather than a retention one, and is
the reason the non-capture classes must not be pooled.

\clearpage

\section{Endogenous exit capacity}
\label{sec:capacity}

The endogenous exit-capacity module replaces the global closure parameter with a
per-agent capacity built from outside social ties and economic independence.
This section reports the closure threshold of the tie channel, the generational
arms, and the anti-circularity test. Model:
\path{src/religion_fundamentalism_abm_v3_3.py}. Aggregated outputs and per-run
artifacts are under \path{results/v3_3_endogenous_capacity_v2/}.

\subsection{Closure of the tie channel}

The tie channel has fixed point
$o^{*} = (\text{renewal base} - \text{crowd-out})/(\text{decay} - \text{renewal})$,
and closes when $o^{*} \leq 0$, which occurs when crowd-out exceeds the base
renewal rate of 0.006. Table~S5.1 reports a sweep over crowd-out at
$\mu = 0.0$ with 5 seeds per level. Source:
\path{results/v3_3_endogenous_capacity_v2/closure_summary.csv}.

\begin{table}[htbp]
\centering
\small
\caption{Table S5.1: Closure sweep over tenure crowd-out.}
\begin{tabular}{lcccc}
\toprule
Crowd-out & Predicted $o^{*}$ & Tie channel closes & Mean final exit capacity & Exit over last 100 steps \\
\midrule
0.0005 & $+0.917$ & no & 0.9302 & 0.0415 \\
0.002 & $+0.667$ & no & 0.8392 & 0.0544 \\
0.008 & $-0.333$ & yes & 0.5166 & 0.0685 \\
\bottomrule
\end{tabular}
\end{table}

The behaviour is a threshold rather than a ratchet. At $\mu=0.0$ the economic
channel remains high, so total capacity floors at 0.5166 when the tie channel
closes rather than falling to zero. Driving capacity to zero requires closing
the economic channel as well, which needs $\mu$ above its own threshold.

\subsection{Generational arms}

With the economic channel inactive at $\mu=0.0$ and 5 seeds per arm, the
born-inside cohort has lower terminal capacity. Source:
\path{results/v3_3_endogenous_capacity_v2/generational_summary.csv}.

\begin{table}[htbp]
\centering
\small
\caption{Table S5.2: Generational arms, economic channel inactive.}
\begin{tabular}{lcc}
\toprule
Born-inside fraction & Exit capacity at step 449 & Cumulative exit \\
\midrule
0.0 & 0.8854 & 0.5747 \\
0.8 & 0.7721 & 0.5680 \\
\bottomrule
\end{tabular}
\end{table}

With the economic channel live, the same comparison is run at
$\mu \in \{0.6, 0.8\}$, both crowd-out levels, both born-inside fractions, and
10 seeds. The economic fixed point is $e^{*} = 0.004/(0.01\mu)$, giving 0.667 at
$\mu=0.6$ and 0.500 at $\mu=0.8$, both below 1, and the observed mean economic
independence of 0.52 to 0.74 confirms that the channel is genuinely interior
rather than clipped at its bound. Source:
\path{results/v3_3_endogenous_capacity_v2/live_econ_summary.csv}.

\begin{table}[htbp]
\centering
\small
\caption{Table S5.3: Generational arms, economic channel live.}
\begin{tabular}{lcccc}
\toprule
$\mu$ & Crowd-out & Born inside & Exit capacity at step 449 & Cumulative exit \\
\midrule
0.6 & 0.002 & 0.0 & 0.7473 & 0.5230 \\
0.6 & 0.002 & 0.8 & 0.6015 & 0.5117 \\
0.6 & 0.008 & 0.0 & 0.4760 & 0.4060 \\
0.6 & 0.008 & 0.8 & 0.3542 & 0.2943 \\
0.8 & 0.002 & 0.0 & 0.6686 & 0.5147 \\
0.8 & 0.002 & 0.8 & 0.5398 & 0.4757 \\
0.8 & 0.008 & 0.0 & 0.4008 & 0.3560 \\
0.8 & 0.008 & 0.8 & 0.3006 & 0.2353 \\
\bottomrule
\end{tabular}
\end{table}

A born-inside fraction of 0.8 gives lower terminal capacity than a fraction of
0.0 in every parameter cell and in all 40 of 40 seed-matched comparisons. Mean
cumulative exit is likewise lower in all four parameter cells; at the level of
individual replicates the comparison holds in 32 of 40 cases, and every reversal
or tie occurs at the lower crowd-out level of 0.002. Both structural knobs act
monotonically: higher economic dependence and higher crowd-out each lower
capacity.

\subsection{Anti-circularity}

The concern this test addresses is that exit capacity might be responding to
enforcement, which would make capture follow by construction. The test proceeds
in two stages over a panel of 93,268 transitions in 300 clusters. Source:
\path{results/v3_3_endogenous_capacity_v2/anticircularity_panel.csv}.

First, the capacity update is reconstructed from its declared inputs, namely
ties, economic independence, tenure, and heterogeneity. The maximum residual
across all transitions is $3.331\times10^{-16}$, which is machine zero. The
update therefore contains no undeclared term.

Second, that structural residual is regressed on lagged punishment received,
lagged punishment delivered, and enforcer status. Table~S5.4 reports the
coefficients.

\begin{table}[htbp]
\centering
\small
\caption{Table S5.4: Regression of the structural residual on enforcement history.}
\begin{tabular}{lccc}
\toprule
Regressor & Coefficient & Standard error & $p$ \\
\midrule
Lagged punishment received & $7.96\times10^{-19}$ & $6.04\times10^{-19}$ & 0.188 \\
Lagged punishment delivered & $5.17\times10^{-20}$ & $9.97\times10^{-20}$ & 0.604 \\
Enforcer status & $-4.83\times10^{-19}$ & $6.17\times10^{-19}$ & 0.434 \\
\bottomrule
\end{tabular}
\end{table}

All three coefficients are of order $10^{-19}$ or smaller, that is at the scale
of the machine-zero residual itself, and none approaches conventional
significance. Enforcement has no structural effect on exit capacity, so the
capacity mechanism is an independent input to capture rather than a restatement
of it.

\clearpage

\section{Computational reproducibility}
\label{sec:repro}

\subsection{Committed run outputs}

Every reported quantity derives from one of the committed aggregate files listed
in Table~S6.1. Each row of a \path{sweep_seed_results.csv} file is one run, and
the per-run directories beneath it hold that run's \path{metrics.csv}, with one
row per step, and \path{agent_summary.csv}, with one row per agent.

The parameter set behind these sweeps is provided in the repository, in the file
\path{results/BEST_PARAMS_frozen.json}, and is reproduced in Table~S1.10.

\begin{table}[htbp]
\centering
\footnotesize
\caption{Table S6.1: Committed aggregate outputs.}
\begin{tabular}{p{4.4cm}p{7.4cm}p{2.0cm}}
\toprule
Experiment & Path & Runs \\
\midrule
Shock factorial, off & \path{recon/factorial_shock_off/sweep_seed_results.csv} & 4{,}050 \\
Shock factorial, on & \path{recon/factorial_shock_on/sweep_seed_results.csv} & 4{,}050 \\
Closure decoupling, $\kappa=1.5$ & \path{recon/decoupled_k1.5/sweep_seed_results.csv} & 4{,}050 \\
Closure decoupling, $\kappa=3.0$ & \path{recon/decoupled_k3.0/sweep_seed_results.csv} & 4{,}050 \\
Closure decoupling, $\kappa=6.0$ & \path{recon/decoupled_k6.0/sweep_seed_results.csv} & 4{,}050 \\
Imposed closure & \path{recon/exogenous_delta_fixed/sweep_seed_results.csv} & 1{,}620 \\
Activation boundary, open & \path{recon/boundary_open/sweep_seed_results.csv} & 1{,}890 \\
Activation boundary, sealed & \path{recon/boundary_sealed/sweep_seed_results.csv} & 1{,}890 \\
Privilege ablation & \path{recon/privilege_ablation/ablation_seed_results.csv} & 120 per arm \\
Endogenous exit capacity & \path{results/v3_3_endogenous_capacity_v2/} & see Section~S5 \\
\bottomrule
\end{tabular}
\end{table}

\subsection{Drivers and aggregators}

Table~S6.2 names the script that produced each set of outputs and the script
that reduces them. The classifier in \path{src/regime_classifier.py} is the
single source of truth for discrete outcomes and carries a self-check.

\begin{table}[htbp]
\centering
\small
\caption{Table S6.2: Drivers and aggregators.}
\begin{tabular}{ll}
\toprule
Role & Script \\
\midrule
Main sweep driver & \path{scripts/run_v2_7_endogenous_delta_sweep.py} \\
Factorial launcher & \path{recon/run_factorial.sh} \\
Decoupling launcher & \path{recon/run_decoupled.sh} \\
Privilege ablation driver & \path{recon/run_privilege_ablation.py} \\
Endogenous capacity driver & \path{scripts/run_v3_3_experiments.py} \\
Boundary aggregator & \path{recon/analyze_boundary.py} \\
Imposed-closure aggregator & \path{recon/analyze_exogenous_fixed.py} \\
Decoupling aggregator & \path{recon/analyze_decoupled.py} \\
Endogenous capacity aggregator & \path{scripts/aggregate_v3_3.py} \\
Outcome classifier & \path{src/regime_classifier.py} \\
\bottomrule
\end{tabular}
\end{table}

\subsection{Per-run configuration records}

Each run of the endogenous exit-capacity experiments writes a
\path{run_config.json} into its own output directory. That file records the full
command line, the resolved model parameters, the creation and completion
timestamps, the model file path, and the SHA-256 digest of the model source
under the key \texttt{code\_sha256}. A reader can therefore confirm that a given
committed run was produced by the exact source file now in the repository.

The main sweeps under \path{recon/} do not write a per-run configuration record.
Their launch parameters are recorded in three places. The swept parameters
appear as columns of \path{sweep_seed_results.csv} and as the directory names
beneath it. The population-level settings \texttt{base\_opp},
\texttt{punish\_floor}, \texttt{exit\_threshold}, and
\texttt{capture\_exit\_cap} appear in the \path{sweep_report.txt} written
alongside it. The seven calibration-supplied parameters appear in neither, and
are instead provided in the repository as the file
\path{results/BEST_PARAMS_frozen.json}, whose values are listed in Table~S1.10
and whose empirical confirmation against the committed outputs is described in
the note to Section~S1.

The absence of a per-run configuration record under \path{recon/} is a
reproducibility weakness of these sweeps rather than of the model. It is the
reason the seven values had to be recovered from run output and verified by
replay, and it is the reason the parameter set is now committed as a file in its
own right. The endogenous exit-capacity runs of Section~S5 do not have this
weakness, because each writes a \path{run_config.json}.

\subsection{Recomputation of outcome classes}

All reported regimes are recomputed from the continuous columns
\texttt{final\_exit\_rate} and \texttt{max\_punish} using the canonical
classifier. No reported quantity reads a stored \texttt{regime} column. This
rule is load-bearing rather than stylistic. The \texttt{regime} column committed
alongside the \path{recon/} sweeps was written under an earlier classifier that
applied a prevalence gate, and for the shock factorial that stored column
disagrees with the canonical classifier on 71 percent of runs, labelling 2,880
runs MIXED that the canonical classifier labels CAPTURE. Stored regime columns
are retained in the committed files as immutable simulation outputs, but they
are not authoritative, and analysis recomputes classification from the
continuous quantities every time.

\end{document}
